\section{Discussion and conclusion}\label{sec:conclusion}
The random-sign IBLT lets one table simultaneously perform decoding and difference-size estimation, with the cells and the decoding capability both unchanged. Whether decoding succeeds is determined only by the mapping, and the randomness of the estimate comes from the signs; this separation makes the estimate exactly unbiased for every fixed mapping and identically distributed for all sign compositions, and makes the post-failure distribution fully characterizable: Gaussian when $d\gg\sqrt M$, and a discrete distribution of a $\pm2$ term plus a Skellam term when $c=d/\sqrt M$ is bounded. The same family of distributions describes all scales, and the Skellam rule that takes it as the null distribution makes the conditional miss rate asymptotically at most the target risk uniformly over all $d\ge2$.

At deployment we recommend combining the two rules by the first-round observation, which keeps the uniform guarantee for all $d\ge2$ with the time in the overload region the same as the Gaussian rule; the guarantees and costs of the various choices are in Section~\ref{sec:exp-deploy}. On the deployment parameters, the Skellam rule meets the target on all experimental configurations with $M\ge256$, and for engineering use we recommend $M\ge1024$ (Section~\ref{sec:exp-cover}). The checksum can be shortened to 32 bits in exchange for narrower cells, leaving the risk of sizing unaffected, and the probability of a verification error is still at most $10^{-4}$ at $M_1=16384$ and load at most 3 (Section~\ref{sec:protocol}). Storing the counts of several independent signs per cell reduces the absolute error but does not change the post-failure conditional scale, and one still performs inversion by the post-failure distribution (Appendix~\ref{app:multi}). The random-sign encoding does not increase the space overhead of the cells, and the increase in decoding time is less than 1\% of the total protocol time (Section~\ref{sec:exp-cost}).

Three questions remain for future work. First, before building the tables, reject more candidates with a rigorous finite-parameter tail upper bound, reducing the computational cost of the Skellam rule. Second, turn the errors of Theorem~\ref{thm:cross} and Theorem~\ref{thm:clt} into computable constants, giving a finite risk bound for a specified $M$. Third, extend to multiset differences whose weights can exceed one: there the random-sign count gives an unbiased estimate of the $F_2$ of the difference vector, and the decoding and post-failure analysis of this extension remain to be completed.
